\chapter{The balance sheet}
\label{ch:ledgerch}

This chapter does not narrate. The previous ten have done that. This one audits:
for each domain, what sixty-seven years produced, stated once, with both sides,
and the causal question asked rather than assumed. It is the hinge of the book,
because it is where the record turns into the constraint set that Part III has
to design within.

\section{Health}

\emph{What the record shows.} Cuba took a country with the best health
indicators in Latin America and a rural half with none of them, closed the
rural--urban gap, and held the national outcomes for sixty years through two
economic collapses. Life expectancy went from around 62 in 1958 to the high
seventies, converging with the United States on a fraction of the spending.
Infant mortality went from around 37 per thousand to single digits and stayed
there --- including through 1990--94, when it continued falling while national
output fell by a third. Polio, measles, diphtheria, rubella and neonatal tetanus
were eliminated. The delivery model, a polyclinic and then a family doctor
responsible for a defined population, is the mechanism, and it has been copied
deliberately by other countries because it works.

\emph{What has happened since 2019.} The system is being hollowed out from three
directions at once. Pharmaceutical supply has collapsed: the domestic industry
cannot obtain inputs and the state cannot pay for imports, and pharmacies are
empty at rates the health ministry itself acknowledges. Physical plant has
degraded past the point where outcome indicators can be maintained by staff
effort --- hospitals without reliable water, power or sterilisation. And the
staff are leaving, in the general emigration and through the medical missions.
Cuba's physician numbers remain high on paper because the paper counts doctors
on mission abroad.

\emph{The honest reading.} The achievement is real, was not bought by
measurement games, and is the single most valuable thing in this ledger. It is
also, in 2026, being lost in real time, and it will not survive another five
years of the present trajectory. That is the argument for
Chapter~\ref{ch:welfare} and the reason it is placed early: the asset is
depreciating faster than the plan can be written.

\section{Education}

\emph{What the record shows.} Illiteracy eliminated in a year and kept
eliminated. Universal primary and secondary schooling. Attainment measured by
an outside body two standard deviations above the Latin American mean, with the
poorest Cuban quartile outperforming the regional average --- which is the
number that matters, because it says the system compresses rather than
reproduces advantage. Higher education free and expanded by an order of
magnitude. Women's attainment ahead of men's since the 1980s.

\emph{The cost side, and the erosion.} The curriculum was a political
formation system, entire disciplines were closed or reduced to catechism, and
Cuba consequently has excellent engineers and physicians and very few people
trained to think about how a market, a court or a company works --- which is a
direct constraint on Part III's implementability, not a matter of taste. And the
system is eroding now: teachers are leaving for the private sector and for
abroad, enrolment in teacher training has fallen, and an informal private
tutoring economy has appeared to fill the gap, which is the classic first stage
of a two-tier system.

\emph{The honest reading.} The best asset Cuba has, the one Part III's second
track and biotechnology sector are both built out of, and the one whose
depreciation is least visible year to year, because a school system's decline
shows up in outcomes a decade later.

\section{Food}

\emph{What the record shows.} A country with good soil, adequate water, a
year-round growing season and a trained agronomic profession imports something
in the range of two thirds to four fifths of what it eats, and cannot pay for
it.\unv{} A substantial share of arable land is idle or under \emph{marab\'u}
scrub. The sugar industry that was three quarters of exports as recently as 1989
produced a harvest in the low hundreds of thousands of tonnes in the mid-2020s
and the country imports sugar.

\emph{The causes, in order of size.} State procurement monopoly setting prices
below cost; tenure that cannot be sold, mortgaged or reliably inherited; no
input market; no functioning transport or cold chain; and a currency in which
nobody wants to be paid. Every one of these is domestic and legal.

\emph{And the money went elsewhere.} In 2024, tourism-related construction took
37.4 per cent of national investment and agriculture 2.7 per cent --- a ratio of
fourteen to one, in a country asking the World Food Programme for children's
milk. Whatever else is disputed about Cuban economic policy, this allocation was
chosen.

\emph{The honest reading.} This is the clearest single case in the book where
the embargo explains very little. Cuba may buy food from anyone in the world and
does buy it from the United States for cash. The reason it does not grow its own
is the design of the domestic agricultural system, which
Chapter~\ref{ch:land} takes apart.

\section{Energy}

\emph{What the record shows.} Sixty years of electricity generation from
imported hydrocarbons supplied on political terms --- Soviet, then Venezuelan
--- burnt in large central thermal plants, most of which are now past their
design lives. When the political supply faltered in 1990 the country went dark;
when it faltered again after 2019 the country went dark, and in October 2024 the
grid failed as a system, nationally, for the first time.

\emph{The honest reading.} Energy is where Cuba's three external dependencies
converge, and it is the binding constraint on everything else --- water,
refrigeration, food processing, health, tourism, industry, connectivity. It is
also the one domain where the technology has changed enough since 2010 that the
answer available now is not the answer that was available then.
Chapter~\ref{ch:energy} accordingly comes before every sector chapter in
Part III.

\section{Demography}

\emph{What the record shows.} Fertility fell below replacement in 1978 and has
never returned; by the mid-2020s it was around 1.2 to 1.4, among the lowest in
the world.\unv{} More than a quarter of the population is over sixty. And
between one and two million people, disproportionately of working and
child-bearing age, left between 2021 and 2025 from a population of eleven
million.

\emph{The honest reading.} This is the hardest constraint in the book and the
least reversible. There is no cohort in the pipeline, so no policy adopted in
2026 changes the working-age population before the 2040s except by migration.
Every proposal in Part III is therefore labour-constrained rather than
capital-constrained, and a pension system at European replacement rates is not
financeable on this dependency ratio at any plausible tax rate --- which is why
Chapter~\ref{ch:welfare} has to say so rather than assume it away.

\section{Rights}

\emph{What the record shows, without inflation and without excuse.} No
independent press since 1961. No legal political organisation other than one. No
independent trade union. No court that has ever ruled against the executive on a
political matter. Summary executions in 1959 with no published count. Labour
camps for believers, homosexuals and nonconformists, 1965--68. A cultural purge
by administrative reassignment, 1971--76. Emigration punished by confiscation
and forced labour for two decades, and mobs organised against departing
neighbours in 1980. Seventy-five people sentenced to up to twenty-eight years in
2003 and three shot in the same month. Thirty-seven people drowned off Havana in
1994 with no independent investigation. Hundreds prosecuted after 11 July 2021,
with sentences of four to thirty years, including of minors.

\emph{What should also be said.} The scale is not that of the region's worst
dictatorships and inflating it is a disservice: Cuba did not run death squads,
did not disappear people in the tens of thousands, and the 20,000-dead figure
attributed to Batista is itself a fabrication that has been recycled for sixty
years by people who should know better. Violent crime is low, the police do not
kill people at Latin American rates, and Cubans move around their own cities at
night in a way that citizens of most countries in the hemisphere do not. Sixty
years of one-party rule produced a repressive state, not a murderous one, and
the distinction matters when Chapter~\ref{ch:polity} designs the transition,
because the settlement a country needs after Argentina is not the settlement it
needs after this.

\section{The embargo question}

This is the question everything else is argued through, and the honest answer
has two halves that are usually published separately.

\emph{What the embargo does cost.} It is not nothing and the dismissive version
is wrong. Shipping: a vessel that calls at a Cuban port cannot enter a United
States port for 180 days, which removes Cuba from the efficient Caribbean
routings and raises freight on everything. Finance: the extraterritorial
enforcement --- of which the multi-billion-dollar penalties imposed on
European banks are the demonstration --- means international banks refuse Cuban
business as a matter of policy rather than of law, so a Cuban transaction is
expensive, slow, or impossible even where it is legal. Medical and industrial
goods: any product with more than a de minimis share of American content is
caught, which in practice covers much modern medical equipment. Credit: Cuba is
excluded from the IMF, the World Bank and the Inter-American Development Bank,
which is not merely a loss of lending but a loss of the technical assistance and
the seal of approval that other countries use to attract private capital. And
Helms--Burton converted all of it from an executive instrument into a statute,
so it cannot be traded away by a president.

\emph{What the embargo does not explain.} Cuba trades freely with the rest of
the world. The European Union, China, Canada, Brazil, Russia, Mexico, Spain,
Vietnam and Venezuela have not embargoed anything, and Cuba has bought from all
of them. The binding constraint on Cuban imports is not permission; it is that
Cuba cannot pay, and cannot pay because it has been in default since 1986, froze
foreign firms' accounts in 2009, and has arrears with almost every supplier who
has extended it credit. That is a credit history and it has no foreign author.
Nor does the embargo explain the domestic prohibitions: the abolition of small
business in 1968, the refusal of legal personality to private firms until 2021,
the absence of a wholesale market, the state procurement monopoly in
agriculture, the ban on Cubans entering hotels until 2008, the exit permit until
2013, the licensing of activity by enumerated list, or the employment agency
that takes the currency spread from workers in foreign ventures. Every one of
those was written in Havana, and every one of them can be repealed in Havana.

\emph{How to hold both.} The formulation this book uses is that the embargo
sets the cost of Cuban failure and the Cuban system sets its probability. A
well-designed Cuban economy under embargo would be poorer than the same economy
without one --- meaningfully poorer, and the difference is measured in
percentage points of growth a year, not in the fantastical cumulative totals
either side publishes. But it would not be an economy that cannot light its
cities or feed itself, because several countries in the region manage both with
worse endowments and no external assistance. The reason this book's Part III
concentrates on domestic institutions is not that the embargo is unimportant.
It is that the embargo is the one variable Cubans do not control, and a design
that waits on it is not a design.

\section{The pattern}

Across sixty-seven years the same sequence occurs four times, and naming it is
the finding Part III has to be built against.

\emph{Duress produces reform.} 1980, 1993--95, 2008--2013, 2021: in each case a
material emergency forced the state to legalise something it had prohibited ---
farmers' markets, dollars and self-employment, usufruct land and private house
sales, private firms with legal personality.

\emph{The reform works.} Food appears; the currency stabilises; the private
sector grows; the shops restock. In every instance the measured result was
positive and visible within months.

\emph{The pressure lifts, and the reform is withdrawn.} 1986, 1996--2004, 2016--17,
2024--25: markets closed, licences frozen, taxes raised, categories deleted,
price caps imposed, and the reformers criticised for concentration of wealth.

\emph{The withdrawal arrives just before it is least affordable.} Rectification
was launched in 1986 and the Soviet Union collapsed in 1990. The licence freeze
came in August 2017, fourteen months before the American re-tightening. The
attack on the new private importers began in 2024, in the middle of the worst
shortages since 1993.

The reason for the pattern is not economic and pretending otherwise wastes
everyone's time. Each reform creates people whose income does not come from the
state, and that is the thing the system is organised to prevent. The reforms are
therefore always framed as concessions, always revocable, and always revoked ---
and because they are known to be revocable, nobody invests behind them, which
guarantees that they underperform, which supplies the evidence for revoking
them.

That is the finding, and it converts directly into the design requirement that
opens Part III: \emph{the binding problem in Cuba is not which policies to
adopt but how to make them irrevocable}. Every measure this book proposes has
been tried in Cuba in some form. Most of them worked. All of them were reversed.
A design whose only content is a better list of policies has not understood the
last thirty years, and Chapter~\ref{ch:polity} is placed second in Part III for
that reason: without a credible commitment mechanism, nothing that follows it is
worth writing down.

\begin{ledger}{1959--2026}
\emph{Achieved.} A health system that closed the rural--urban gap, converged
with rich-country outcomes on poor-country money, and held its indicators
through a 35 per cent contraction. An education system whose measured results
beat its region by two standard deviations and compress rather than reproduce
advantage. Illiteracy ended in a year. Legal and enforced racial desegregation.
Universal social security, near-universal home ownership, and one of the lowest
violent crime rates in the hemisphere. A biotechnology sector that developed and
deployed its own vaccines when far richer countries could not. A decisive
contribution to ending South African rule in Namibia. And a state that has never
lost the capacity to reach every block of its territory, which is why it
absorbed 1993 without famine.

\emph{Cost.} An export basket unchanged in structure from 1958 to 1989 and
destroyed since. Three successive external patrons and three identical
failures to diversify under them. Sugar below nineteenth-century output and
food imported at two thirds of consumption in a fertile country. A grid that
failed as a system in 2024. A default unresolved since 1986. Six decades
without a press, a union, a party or a court that was not the state's. Camps,
purges, exile as punishment, mobs organised against neighbours, and
thirty-year sentences in 2021 for walking in the street. Fertility below
replacement since 1978 and between one and two million people gone in five
years. And a reform record in which everything that worked was later withdrawn,
which is the reason none of it compounded.
\end{ledger}
